\section{A finite tangent-field theorem}\label{sec:table}

Put $d=n+1$. Realize the boundary of the $d$-dimensional
cross-polytope as $P=\{p:\|p\|_1=1\}$, with vertices
$+e_1,-e_1,\ldots,+e_d,-e_d$. We use only the rational vertices
of its $r$-fold barycentric subdivision, denoted here by $K_r$.
Its mesh is the largest $\ell^1$ distance between two vertices of one
simplex. The norms in this section are the rational operations
$\|a\|_1=\sum_i|a_i|$ and $\|a\|_\infty=\max_i|a_i|$.

\begin{theorem}[Finite small-vector witness]\label{thm:finite-table}
Suppose $n$ is even, $\tau>0$ is rational, and the mesh of $K_r$ is
at most $1/(8d)$. At every vertex $p$ let a rational vector $v_p$
be given, with $p\cdot v_p=0$. Assume that for every edge $[p,q]$,
\[
             \|v_p-v_q\|_1\leq \frac{\tau}{16d}.
\]
Then a bounded scan of the vertices returns a vertex $p$ with
$\|v_p\|_1\leq\tau$.
\end{theorem}

The geometry of the failed-scan argument takes place in a plane. If a
vertex value is large, rescale it to a tangent vector $w_p$ of unit
$\ell^1$ norm. Join $p$ to $w_p$, then $w_p$ to $-p$, by straight
segments. Orthogonality keeps both segments away from the origin.
The proof measures that separation uniformly over all vertices.
\begin{center}
\begin{tikzpicture}[x=1cm,y=1cm,>=stealth]
  \draw[gray!65] (-2.1,0)--(2.1,0);
  \draw[gray!65] (0,0)--(0,1.03);
  \draw[gray] (0.15,0)--(0.15,.15)--(0,.15);
  \draw[ink,thick,->] (1.8,0)--(.06,1.03);
  \draw[ink,thick,->] (0,1.06)--(-1.74,.04);
  \fill[ink] (1.8,0) circle (1.3pt) node[below] {$p$};
  \fill[ink] (-1.8,0) circle (1.3pt) node[below] {$-p$};
  \fill[ink] (0,1.06) circle (1.3pt) node[above] {$w_p$};
  \fill (0,0) circle (1.1pt) node[below] {$0$};
\end{tikzpicture}
\end{center}

\begin{proof}
For $n=0$, the two directions are $p=\pm1$ and tangency forces
$v_p=0$. For $n>0$, scan the table in vertex order and return the
first entry satisfying the conclusion. The remaining branch has
$A_p=\|v_p\|_1>\tau$ at every vertex. Define $w_p=v_p/A_p$.
Then $p\cdot w_p=0$ and $\|p\|_1=\|w_p\|_1=1$. If $p,q$
are vertices of one simplex, the edge hypothesis and the reverse
triangle inequality give
\[
 \|w_p-w_q\|_1
 \leq \frac{\|v_p-v_q\|_1+|A_q-A_p|}{A_p}
 \leq \frac{2\|v_p-v_q\|_1}{\tau}
 \leq \frac1{8d}.
\]
This includes $p=q$.

Take $L=16d$. Make $2L+1$ tables $H_i$, each with one rational
vector per vertex, using the two formulas
\[
 \begin{split}
 H_j(p)&=(1-s)p+s w_p,\\
 H_{L+j}(p)&=(1-s)w_p-s p,
 \qquad s=j/L,\quad 0\leq j\leq L .
 \end{split}
\]
The two formulas agree at $H_L$. The squared length of any entry
is at least $((1-s)^2+s^2)/d$. Indeed, the cross term vanishes by
tangency, and for any $a$ of $\ell^1$ norm one the identity
\[
 d\sum_i a_i^2-\Big(\sum_i|a_i|\Big)^2
   =\sum_{i<j}(|a_i|-|a_j|)^2
\]
gives $\sum_i a_i^2\geq 1/d$. As $(1-s)^2+s^2\geq 1/2$,
and $\sum_j (H_i(p)_j)^2\leq d\|H_i(p)\|_\infty^2$, it follows
by rational comparison that
\[
                      \|H_i(p)\|_\infty\geq\frac{1}{2d}.
\]
At any one time the $\ell^1$ distance between two entries on the
same simplex is at most $1/(8d)$, by the mesh bound and the estimate
on $w_p-w_q$. At a fixed vertex one time step changes the entry by
at most $2/L$ in $\ell^1$, since each endpoint vector of each leg has
norm one. Consequently, for $p,q$ in a common simplex and
$|i-i'|\leq1$,
\[
       \|H_i(p)-H_{i'}(q)\|_\infty
        \leq \frac{1}{8d}+\frac{2}{L}=\frac{1}{4d}.
\]

Label $p$ at time $i$ by the signed basis vector corresponding to the
first coordinate of $H_i(p)$ of largest absolute value, with its
observed sign. If entries on the same simplex at equal or adjacent
times received labels $+e_j$ and $-e_j$, their $j$th coordinates
would differ by at least $1/d$. This contradicts the preceding upper
bound. A collection of signed basis vertices containing no opposite
pair spans a face of the cross-polytope. Hence the labels at each
time define a simplicial map, and maps at consecutive times are
contiguous because the union of their labels on each source simplex
spans a target face.

At time zero the label $\lambda(p)$ is the signed largest coordinate
of $p$ with the same tie convention. Its coordinate is nonzero, so
it belongs to the smallest original face containing $p$. Thus
$\lambda$ is a carrier map. At the final time the labels are
$-\lambda(p)$, because the first maximizing index is unchanged
by negation. The contiguous sequence from $\lambda$ to $-\lambda$
contradicts Theorem~\ref{thm:obstruction}. This eliminates the
unsuccessful branch of the bounded scan.
\end{proof}

The computation returning $p$ only compares the entries of the
original table with $\tau$. Normalized vectors and labels certify
why this scan succeeds. Their denominators are positive on the branch
where they are used, and the proof uses only rational arithmetic and
bounded loops.
